When a quadruped snags a leg on a wire, a table leg, or loose netting, the fault is easy to miss with external cameras or lidar, and the proprioceptive trace it leaves is faint and easily mistaken for ordinary standing or for the rigid post-stand-up ``Lock'' stance. We treat the problem as time-series event detection and train a single multi-task causal temporal convolutional network on a $0.40$~s, $60$-channel window of joint, foot-contact, and inertial signals. The network jointly predicts whether a leg is entangled, which of the four legs is affected, and a per-leg severity in $[0,1]$ that is grounded in a physics index combining a confidence gate, a Mahalanobis deviation from the walking distribution, and resistive thigh effort. Because every output is read from the last timestep of a strictly left-padded (causal) stack, the model validated offline is the one deployed through an ONNX runtime, whose outputs match the research pipeline to within $4\times10^{-7}$, so the architecture introduces no train--deploy gap. On a leakage-safe held-out split grouped by recording, the detector reaches an F1 of $0.809$, a ROC-AUC of $0.956$, and a multi-label exact-match of $0.908$; leave-one-recording-out cross-validation over $15$ folds gives an F1 of $0.807 \pm 0.142$. Inference costs about $1$~ms per window on a single CPU thread. All results are from recorded-log replay; we report offline evaluation only and make no claim of live on-robot detection performance.
